%% ============================================================
%% Abstract (English)
%% ============================================================
\addcontentsline{toc}{chapter}{Abstract}
\chapter*{Abstract}

Mandarin Chinese is a tonal language in which syllable-level pitch patterns distinguish lexical meaning, making tone production a central challenge for learners whose first language is non-tonal. This thesis investigates automatic Mandarin tone classification for Polish L1 learners as the assessment component of a future computer-assisted pronunciation training system. Pronunciation quality is decomposed into two largely independent dimensions, phonetic accuracy of the consonant--vowel realisation and tonal accuracy of the produced tone; the phonetic component is addressed in a parallel project on the same learner corpus, while this thesis develops and evaluates the tonal-accuracy classifier. Six modelling setups are evaluated: a ToneNet-style convolutional neural network (CNN) baseline using mel-spectrogram images, three Wav2Vec 2.0 variants with different adaptation routes, and two Chinese HuBERT-Large variants with different prediction heads.

The experiments use an expert-labelled learner corpus of Mandarin recordings produced by Polish speakers. The self-supervised learning (SSL) models are evaluated on a speaker-disjoint split containing 8,577 training utterances, 464 development utterances, and 2,282 test utterances, while the CNN baseline follows the original ToneNet-style random-split protocol. Because the thesis compares complete experimental setups, the split protocol is treated as part of each setup definition and is reported explicitly. Across the evaluated setups, Chinese HuBERT-Large with frame-level cross-entropy achieves the highest test accuracy, macro-F1, and weighted F1, reaching 79.93\% accuracy, 73.85\% macro-F1, and 79.89\% weighted F1. Among the Wav2Vec 2.0 routes, the strongest macro-F1 is obtained by the AISHELL-1-initialised label2 route, reaching 72.06\%. These results suggest that frame-level supervision over SSL representations is effective for learner-tone classification, while T3 and non-standard tonal productions remain the most difficult categories.

\section*{Keywords:}
\noindent Mandarin tone classification, computer-assisted pronunciation training, self-supervised learning, Wav2Vec 2.0, HuBERT, Polish L1 learners
